\item De manera subconsciente, el ser humano estima la distancia a un objeto desde el ángulo que subtiende su campo visual. Este ángulo $\theta$ en radianes está relacionado con la altura lineal $h$ del objeto y la distancia $d$, de acuerdo con $\theta = h/d$. Suponga que está manejando un automóvil y que otro vehículo, de $1.50\text{ m}$ de altura, está a $24.0\text{ m}$ detrás de usted.
\begin{enumerate}
    \item Suponga que su automóvil tiene un espejo retrovisor de tipo plano del lado del pasajero, a $1.55\text{ m}$ de los ojos del conductor. ¿A qué distancia de los ojos del conductor está la imagen del automóvil que le viene siguiendo?
    \item ¿Cuál es el ángulo subtendido por la imagen que aparece en su campo visual?
    \item \textbf{¿Qué pasaría si?} Ahora suponga que su automóvil tiene un espejo retrovisor convexo con un radio de curvatura de $2.00\text{ m}$ (figura 35.15). ¿Qué tan lejos de sus ojos está la imagen del automóvil que viene atrás?
    \item ¿Cuál es el ángulo que subtiende la imagen en sus ojos?
    \item En términos de su tamaño angular, ¿qué tan lejos parece estar el automóvil que lo viene siguiendo en comparación con un espejo plano?
\end{enumerate}